% AUTHORED, NOT EXTRACTED: the v1.1 manuscript contains no propositional
% Hilbert calculus.  Canonical source: Mendelson, Introduction to Mathematical
% Logic (1997), axiom schemas A1-A3 with modus ponens (after Lukasiewicz);
% cf. Kleene, Introduction to Metamathematics (1952).  The negation-primitive
% third schema is chosen to match the book's connective conventions, in which
% negation is primitive in the active classical language.
\calculussheet{sheet:hilbert-cpl}{Hilbert Calculus for Classical Propositional Logic}{Language: the $\neg/\to$ fragment of $\mathsf{Fm}_{\mathsf{CPL}}$ (metavariables $\varphi,\psi,\chi$ range over formulas of that fragment; the book's remaining primitives $\bot,\land,\lor$ are reached by the standard classical definitions when this calculus is used for the full language); sequent shape: no sequents --- Hilbert-style derivability $\Gamma\vdash\varphi$ via finite sequences of formulas, each a member of $\Gamma$, an instance of an axiom schema, or obtained from earlier lines by modus ponens; negation: primitive (the third schema governs it directly); equality: none (propositional language).}
\hbrule{rule:hilbert-cpl-a1}{A1}{}{\varphi\to(\psi\to\varphi)}{\none}
\hbrule{rule:hilbert-cpl-a2}{A2}{}{(\varphi\to(\psi\to\chi))\to((\varphi\to\psi)\to(\varphi\to\chi))}{\none}
\hbrule{rule:hilbert-cpl-a3}{A3}{}{(\neg\psi\to\neg\varphi)\to((\neg\psi\to\varphi)\to\psi)}{\none}
\hbrule{rule:hilbert-cpl-mp}{MP}{\varphi\to\psi\qquad\varphi}{\psi}{\none}
